
\subsubsection{Documentation Frameworks}

Datasheets for Datasets \citep{Gebru2018} established lifecycle categories through iterative refinement with dataset creators at Microsoft, Google, and IBM. The framework includes seven sections: motivation, composition, collection process, preprocessing, uses, distribution, and maintenance. Roman et al. (2023) condensed this to a two-tier structure (General Information versus Responsible AI) validated through formative evaluations with dataset producers on HuggingFace. Their Open Datasheets system automates metadata extraction within HuggingFace using structured templates.

Similar frameworks include Model Cards \citep{Mitchell2019} for model documentation and Data Statements \citep{Bender2018} for linguistic datasets. These approaches impose structure through templates filled manually or via guided interfaces, succeeding within repositories that adopt the standards but not addressing cross-repository semantic mapping for heterogeneous metadata.

\subsubsection{Metadata Assessment Tools}

FAIR principles \citep{Wilkinson2016} motivated automated metadata quality assessment. F-UJI \citep{Devaraju2021} and FAIRshake \citep{Clarke2019} evaluate metadata completeness within single repositories using repository-specific rules. Cross-repository comparison requires manual field mapping. Schema.org and Dublin Core provide standardized vocabularies but require explicit adoption, which major repositories implement differently or incompletely.

Semantic embeddings offer a potential approach to cross-repository mapping without explicit ontologies. Sentence-BERT \citep{Reimers2019} captures distributional similarity in frozen representations. This study tested whether such embeddings encode lifecycle categories as geometric structure accessible to unsupervised clustering.

\subsubsection{Clustering Methods}

K-means clustering \citep{Lloyd1982Kmeans} assumes balanced, well-separated clusters in Euclidean space. Topic models like Latent Dirichlet Allocation \citep{Blei2003} discover latent structure through co-occurrence patterns. These methods succeed when distributional similarity aligns with semantic categories and classes are reasonably balanced. Imbalanced clustering methods exist (Chawla et al., 2002; He \& Garcia, 2009) but primarily address supervised classification rather than unsupervised discovery with extreme imbalance ratios.

This study contributes an empirical demonstration that semantic embeddings encode lifecycle roles with high supervised accuracy (79.6\%) yet fail unsupervised recovery (NMI 0.023), identifying class imbalance (11:1 ratio) as a barrier for standard clustering algorithms.

